\section{Data and Labels}\label{sec:data}
\subsection{Video corpus and cue-verified labels}
We collected 24 instructor-led yoga videos (12 original and 12 added later), 1{,}206{,}391 frames in total, and ran MediaPipe on every frame. Frame labels are \emph{cue-verified}: a frame is a positive example of a pose only if (i)~the instructor names that pose in the speech transcript (Whisper-large-v3 \cite{radford2023whisper}, English and Hindi), \emph{and} (ii)~the app's own rule engine agrees with the name, \emph{and} (iii)~the body is held steady. This yields about 80\,k positive frames; about 421\,k frames are labelled transition (the body is moving between poses) and about 705\,k are excluded as ambiguous. The point of the triple check is that no one of the three sources is trusted alone: an instructor may name a pose before or after the body reaches it, and labels from a rule engine alone would only teach a model to copy the rules.

\subsection{Features, and a train/serve mismatch}\label{sec:features}
The models see 15 joint angles computed from the 2-D landmarks (the $z$ coordinate is set to zero): left/right elbow, shoulder, hip, knee and ankle angles; left/right trunk angles; the neck angle; and left/right hip abduction. This is exactly what the browser can compute cheaply and what is sent. An audit showed that the original training table had been built with \emph{raw} 3-D angles while the app serves \emph{zero-$z$} angles: the median per-feature gap was $14^\circ$, a pure train/serve skew \cite{sculley2015debt}. All training tables used here are rebuilt with the served recipe.

\subsection{Photo sets}
Photos supply the variety that 24 videos cannot. \emph{Commons-422} is a set of 422 freely licensed Wikimedia photos with a frozen test split of 103 (hash-split, fixed since 19 Sep 2026). \emph{Public photos} are 6{,}616 photos from two public datasets: 2{,}634 labelled with our poses and 3{,}982 of poses outside our vocabulary, used as negatives; 4{,}931 are used for training and \textbf{1{,}685 are held out} (1{,}037 of them are ``other poses''). Table~\ref{tab:data} summarises the data.
\begin{table}[t]
\centering\footnotesize
\caption{Data used in this paper.}\label{tab:data}
\begin{tabularx}{\columnwidth}{@{}lX@{}}
\toprule
Item & Size / detail \\ \midrule
Training videos & 24, 1{,}206{,}391 frames; $\sim$80\,k positive / $\sim$421\,k transition / $\sim$705\,k excluded \\
Features & 15 zero-$z$ joint angles per frame (MLP); 60$\times$33$\times$3 landmark windows (ST-GCN) \\
Commons-422 & 422 photos; frozen test split 103 \\
Public photos & 6{,}616 (4{,}931 train / \textbf{1{,}685 held out}, of which 1{,}037 are ``other poses'') \\
ST-GCN windows & 24 videos, 8{,}002 windows (60 frames, stride 12), 3 folds \emph{by video} \\ \bottomrule
\end{tabularx}
\end{table}

\subsection{Evaluation protocol}
A number counts only if it was measured on data the model never trained on. A pose \emph{passes} if recall $\ge0.70$ \emph{and} precision $\ge0.70$ with at least 10 held-out photos (5 on the tiny frozen set; 20 windows for the ST-GCN). A \emph{false alarm} is a photo of a pose that is not one of the app's poses that the system names as one of them (lower is better). Sequence models are scored on three folds split \emph{by video}, so the model never sees the person or room it is tested on. The trainers' own validation accuracy (90--98\%) comes from a random split and is never quoted as a result.
